%%%%%%
%
% $Author: Sadegh Naderi $
% $Datum: 2023-10-15  $
% $Pfad: ML23-01-Keyword-Spotting-with-an-Arduino-Nano-33-BLE-Sense\Presentations\SadeghNaderi\LiteratureReview\slides\ContentLiteratureReview.tex $
% $Version: 3.0 $
% $Review Date: 2023-12-16 $
%
%
% !TeX encoding = utf8
% !TeX root = Rename
%
%%%%%%


\Mysection{Automatic Speed Control of DC Series Motor by Arduino  }

\STANDARD{Description}
{ 
	 Der Artikel  \cite{Ali:2023} beschäftigt sich mit Gleichstrom-Reihenschlussmotoren und deren Eigenschaften, insbesondere dem hohen Anlaufdrehmoment und dem veränderlichen Drehzahl-Drehmoment-Verhalten. Ein Schwerpunkt liegt auf verschiedenen Methoden zur Drehzahlregelung, wobei die Steuerung über die Klemmenspannung hervorgehoben wird.r 
	
	\bigskip
	
	\textbf{Abgrenzung:} während der Artikel sich mit der automatischen Geschwindigkeitsregulierung von Gleichstrommotoren beschäftig, geht er wenig auf die eigentlichen Regelungsalorithmen und Methoden ein. Er beschäftigt sich vorwiegend mit Analogen Methoden. 
	
	\scriptsize{\textbf{Keywords:} Motor, Arduino, Speed Control, motor controller, voltag}
}




\Mysection{BLDC Motor Speed Control Based on MPC Sliding
	Mode Multi-Loop Control Strategy – Implementation
	on Matlab and Arduino Software}
\STANDARD{Description}
{ 
	
Dieser Konferenzartikel \cite{Aghaee:2018}  behandelt die Entwicklung, Modellierung und praktische Umsetzung eines fortschrittlichen Regelverfahrens für BLDC-Motoren. Hierbei geht es um einen neurartigen Kaskadenregler aus zwei Regelkreisen ( Slide-Mode Regler für Motorstrom und Model-Predictive-Controller, Drehzahl auf soll Gesteuert), mit dem Zielfokus auf schnelle und stabile Drehzahlregelung, Robustheit gegenüber Laständerungen und Modellunsicherheiten
und der Einhaltung von Stromgrenzen. 
	
	\bigskip
	
	\textbf{Limitation:} Der Artikel erwähnt zwar herkömmliche Methoden zu Drehzahlregulierung, führt jedoch keinen qunatitativen Vergleich durch. Somit werden Potenzial aufgezeigt, jedoch nicht messbar begründet. Auch werden Rechenaufwand für MPC und nichtideale Motoreffekte nicht näher beleuchtet. 
	
	\scriptsize{\textbf{Keywords:} PID, motor control, Matlab, BLDC, control methods,   }
}


\Mysection{PID wihtoud a Phd }
\STANDARD{Description}
{ 
	
	Dieser Artikel \cite{Wescott:2018}	erläutert die grundlegende Funktionsweise der PID Regelung anschaulich anhand eines Elektromotors. Zudem zeigt er Unterschiede einzelner aus der PID abgeleiteter Methoden quantifiziert auf. 
	\bigskip
	
	\textbf{Limitation:} Dieser Artikel stamm einzig von einem Berater für Regelungstechnik und eignet sich somit nur bedingt als Grundlage für einen Bericht. 
	
	\scriptsize{\textbf{Keywords:} PID, PID controll, BLDC, control systems, proportional integral control. }
}



\Mysection{PID-basierte digitale
	Regelungstechnik}
\STANDARD{Description}
{ 
Diese Buch  \cite{Ibrahim:2023} erläutert die PID-basierte Regelungstechnik anschaulich. Es geht vordergründig um die Anwendung von PID Regelung bei Mikrocontrollern wie dem Arduino Uno. Hierbei werden gängigen Komponenten von Regelungsystemen erklärt und anschaulich für Beispiele verwendet. 
	
	\bigskip
	
	\textbf{Limitation:} keine 
	
	\scriptsize{\textbf{Keywords:} Arduino Uno, Mikrocontroller, PID, Regelungstechnik, Sensoren, Elektronik, Kommunikation }
}

\Mysection{ Motorsteuerung mit Microcontollern}
\STANDARD{Description}
{ 
	Diese Arbeit  \cite{Niemeyer:2015} behandelt die Ansteuerung von Motoren mithilfe von Mikrocontrollern. Hierbei wird die grundlegende Funktionsweise der Schnittstellen erläutert und anschließend beispielhaft angewendet. Ebenfalls werden Sensoren und Ihre Relevanz in der Motorsteuerung beleuchtet. 
	
	\textbf{Limitation:} Dies ist eine von der TUM veröffentlichte Studentische Arbeit. Ihre Zuverlässigkeit ist als gering zu betrachten, da lediglich wenige Personen diese Veröffentlichung gelesen bzw. korrigiert haben. 
	
	\scriptsize{\textbf{Keywords:}  Mikrocontroller, Lochscheibe, Regelungstechnik, Sensoren, Elektronik, Kommunikation, Servomotor }
}
\Mysection{PID-basierte digitale
	Regelungstechnik}
\STANDARD{Description}
{ 
	Diese Buch  \cite{Ibrahim:2023} erläutert die PID-basierte Regelungstechnik anschaulich. Ers geht vordergründig um die Anwendung von PID Regelung bei Mikrocontrollern wie dem Arduino Uno. Hierbei werden gängigen Komponenten von Regelungsystemen erklärt und anschaulich für Beispiele verwendet. 
	
	\bigskip
	
	\textbf{Limitation:} keine 
	
	\scriptsize{\textbf{Keywords:} Arduino Uno, Mikrocontroller, PID, Regelungstechnik, Sensoren, Elektronik, Kommunikation }
}